\chapter{\texorpdfstring{$\Ck{k}$}{Ck} maps on sets with dense interior}
\label{chap:denseinterior}

Cubes $[0, 1]^n$, half-spaces and their open subsets are not open, but they are the domains of
the charts of manifolds with boundary or with corners, and of the curves on compact intervals.
Glöckner and Neeb extend the calculus to \emph{locally convex subsets with dense interior}
\cite[\S1.4]{GN26}: a map is $\Ck{1}$ on such a set $U$ if it is $\Ck{1}$ on the interior $U^0$
and its differential extends continuously to $U \times E$. In this chapter we prove that the
first-order calculus of Chapter~\ref{chap:differentiability}, the higher differentials of
Chapter~\ref{chap:higher}, and the $\Ck{k,\ell}$ maps and the exponential law of
Chapter~\ref{chap:mapping} all survive on such sets. Manifolds with rough boundary, whose charts
take values in these sets \cite[Definition~3.5.1]{GN26}, are built on this chapter in
lc-manifolds.

As everywhere else, no completeness is assumed: integrals are weak integrals. In a few places
the proofs depart from those of \cite{GN26}; the remarks say where.

\section{Locally convex subsets and \texorpdfstring{$\Ck{k}$}{Ck} maps}

\begin{definition}[Locally convex subsets with dense interior {\cite[Definitions~1.4.1 and 1.4.4]{GN26}}]
  \label{def:lc-subset}
  \lean{LocallyConvex.Analysis.IsLocallyConvexSubset, LocallyConvex.Analysis.HasDenseInterior,
    LocallyConvex.Analysis.dLC, LocallyConvex.Analysis.ContDiffLC}
  \leanok
  \uses{def:ck}
  A subset $U$ of a locally convex space $E$ is \emph{locally convex} if every point of $U$ has a
  convex neighbourhood in $U$, and it has \emph{dense interior} if $U \subseteq \overline{U^0}$.
  A map $f \colon U \to F$ on such a set is $\Ck{0}$ if it is continuous. It is $\Ck{k+1}$ if it
  is continuous, $f|_{U^0}$ has all directional derivatives, and the differential
  $\dd(f|_{U^0}) \colon U^0 \times E \to F$ extends to a map $\dd f \colon U \times E \to F$ which
  is $\Ck{k}$ on $U \times E$.
\end{definition}

The definition is recursive, as Definition~\ref{def:ck}; Glöckner and Neeb state it through the
iterated differentials, and the two forms agree by Theorem~\ref{thm:lc-higher} below.

\begin{proposition}[Basic properties]
  \label{prop:lc-basic}
  \lean{LocallyConvex.Analysis.contDiffLC_iff_contDiffMBOn,
    LocallyConvex.Analysis.dLC_eqOn_of_continuousOn, LocallyConvex.Analysis.ContDiffLC.of_hasDirDerivAt,
    LocallyConvex.Analysis.ContDiffLC.exists_clm}
  \leanok
  \uses{def:lc-subset}
  Let $U \subseteq E$ be a locally convex subset with dense interior.
  \begin{enumerate}
    \item On open sets, $\Ck{k}$ in the sense of Definition~\ref{def:lc-subset} is $\Ck{k}$ in the
      sense of Definition~\ref{def:ck}.
    \item The continuous extension of $\dd(f|_{U^0})$ is unique.
    \item $f$ is $\Ck{k+1}$ as soon as $\dd(f|_{U^0})$ has some $\Ck{k}$ extension to $U \times E$,
      and then $\dd f$ is that extension.
    \item If $f$ is $\Ck{1}$, then $\dd f(x, \cdot)$ is continuous and linear for every $x \in U$
      \cite[Lemma~1.4.7]{GN26}.
  \end{enumerate}
\end{proposition}

\begin{proof}
  \leanok
  (b) holds because $U^0$ is dense in $U$; (c) is a reformulation of the definition, and all the
  rules below are proved from it. (d): $\dd f(x, \cdot)$ is a pointwise limit of the linear maps
  $\dd f(y, \cdot)$, $y \in U^0$, so it is linear, and it is continuous by definition.
\end{proof}

\section{First-order calculus}

\begin{theorem}[Mean value theorem and difference quotients {\cite[Proposition~1.4.8, Lemma~1.4.9]{GN26}}]
  \label{thm:lc-mean-value}
  \lean{LocallyConvex.Analysis.ContDiffLC.integral_dLC,
    LocallyConvex.Analysis.ContDiffLC.tendsto_quotient}
  \leanok
  \uses{prop:lc-basic, prop:mean-value-weak}
  Let $f \colon U \to F$ be $\Ck{1}$.
  \begin{enumerate}
    \item If the segment $[x, x + h]$ lies in $U$, then
      $\lambda(f(x + h) - f(x)) = \int_0^1 \lambda(\dd f(x + th, h))\, dt$ for every continuous
      linear functional $\lambda$ on $F$.
    \item The difference quotients $t^{-1}(f(x + th) - f(x))$ extend continuously to $t = 0$, with
      value $\dd f(x, h)$.
  \end{enumerate}
\end{theorem}

\begin{proof}
  \leanok
  For a convex $U$ we shrink the segment towards an interior point, where
  Proposition~\ref{prop:mean-value-weak} applies, and pass to the limit in a parameter-dependent
  integral; this gives (a) for convex sets, and (b). For a general locally convex $U$, the
  function $t \mapsto \lambda(f(x + th))$ has right derivative $\lambda(\dd f(x + th, h))$ on
  $[0, 1)$ by (b), and (a) follows from the fundamental theorem of calculus for real functions.
\end{proof}

\begin{remark}
  Glöckner and Neeb prove the mean value theorem first on convex sets, with an integral in a
  completion of $F$, and then split a segment into finitely many pieces lying in convex
  neighbourhoods. The argument above avoids the completion. We do not claim it is new.
\end{remark}

\begin{theorem}[Rules of calculus {\cite[Propositions~1.4.10 and 1.4.11]{GN26}}]
  \label{thm:lc-rules}
  \lean{LocallyConvex.Analysis.ContDiffLC.comp, LocallyConvex.Analysis.ContDiffLC.dLC_comp,
    LocallyConvex.Analysis.ContDiffLC.prodMk, LocallyConvex.Analysis.ContDiffLC.prodMap,
    LocallyConvex.Analysis.ContDiffLC.clm_comp, LocallyConvex.Analysis.ContDiffLC.add,
    LocallyConvex.Analysis.ContDiffLC.bilin, LocallyConvex.Analysis.ContDiffMBOn.contDiffLC,
    LocallyConvex.Analysis.ContDiffLC.mono, LocallyConvex.Analysis.ContDiffLC.of_local,
    LocallyConvex.Analysis.contDiffLC_one_iff_partials,
    LocallyConvex.Analysis.ContDiffLC.dLC_eq_add_partials}
  \leanok
  \uses{thm:lc-mean-value, thm:chain-rule}
  On locally convex subsets with dense interior:
  \begin{enumerate}
    \item \textbf{Chain rule.} If $f \colon U \to V$ and $g \colon V \to G$ are $\Ck{k}$, so is
      $g \circ f$, with $\dd(g \circ f)(x, h) = \dd g(f(x), \dd f(x, h))$.
    \item Pairs, products of maps, compositions with continuous linear maps, sums and continuous
      bilinear images of $\Ck{k}$ maps are $\Ck{k}$; restrictions of $\Ck{k}$ maps on open sets
      are $\Ck{k}$.
    \item $\Ck{k}$ is inherited by locally convex subsets with dense interior, and it is a local
      property.
    \item \textbf{Partial differentials.} A map on $U_1 \times U_2$ is $\Ck{1}$ if and only if its
      partial differentials exist and extend continuously, and then
      $\dd f((x, y), (h, k)) = \dd_1 f((x, y), h) + \dd_2 f((x, y), k)$.
  \end{enumerate}
\end{theorem}

\begin{proof}
  \leanok
  (a) The candidate differential $\dd g(f(x), \dd f(x, h))$ is continuous, and on $U^0$ it is the
  derivative of $g \circ f$ by Theorem~\ref{thm:lc-mean-value}(b) applied to $g$ along
  $f(x + th)$; conclude with Proposition~\ref{prop:lc-basic}(c), by induction on $k$. The chain
  rule needs only $V$ to be locally convex with dense interior, and only density of $U^0$ in $U$.
  (b)--(d) follow from (a) and Proposition~\ref{prop:lc-basic}(c) as in the open case.
\end{proof}

Items (b) and (c) are what the transition maps of manifolds with rough boundary need
\cite[Remark~3.5.2]{GN26}.

\section{Higher differentials, curves and subspaces}

\begin{theorem}[Higher differentials {\cite[Lemma~1.4.5]{GN26}}]
  \label{thm:lc-higher}
  \lean{LocallyConvex.Analysis.iteratedDLC, LocallyConvex.Analysis.ContDiffLC.iteratedDLC_extends,
    LocallyConvex.Analysis.iteratedDLC_unique, LocallyConvex.Analysis.contDiffLC_of_extends,
    LocallyConvex.Analysis.contDiffLC_iff_extends}
  \leanok
  \uses{prop:lc-basic, thm:ck-iff-bastiani}
  A map $f \colon U \to F$ is $\Ck{k}$ if and only if it is continuous, $f|_{U^0}$ is $\Ck{k}$,
  and the iterated differentials $\dd^{(j)}(f|_{U^0})$ extend continuously to
  $U \times E^j$ for $j \le k$. The extensions are unique.
\end{theorem}

\begin{proof}
  \leanok
  Necessity by induction on $k$, from the definition. For sufficiency, \cite{GN26} uses the
  version of \cite[Lemma~1.3.12]{GN26} for sets with dense interior
  \cite[Remark~1.4.12]{GN26}. We argue by induction on $k$ instead, with the explicit formula of
  Theorem~\ref{thm:ck-iff-bastiani} for the differentials of $\dd(f|_{U^0})$ on the open set
  $U^0 \times E$, whose right-hand side extends continuously. We do not claim this argument is
  new.
\end{proof}

\begin{proposition}[Curves and subspaces {\cite[Lemmas~1.4.13 and 1.4.16]{GN26}}]
  \label{prop:lc-curves}
  \lean{LocallyConvex.Analysis.contDiffLC_iff_contDiffCurveOn,
    LocallyConvex.Analysis.ContDiffLC.dLC_curve, LocallyConvex.Analysis.contDiffLC_subtype_iff,
    LocallyConvex.Analysis.IsLocallyConvexSubset.exists_seq_interior}
  \leanok
  \uses{def:lc-subset}
  \begin{enumerate}
    \item On a non-degenerate interval $I \subseteq \R$, the $\Ck{k}$ maps in the sense of
      Definition~\ref{def:lc-subset} are the $\Ck{k}$ curves, with
      $\dd\gamma(s, r) = r\gamma'(s)$.
    \item Let $F_0 \subseteq F$ be a sequentially closed vector subspace. A map $U \to F_0$ is
      $\Ck{k}$ into $F$ if and only if it is $\Ck{k}$ into $F_0$.
  \end{enumerate}
\end{proposition}

\begin{proof}
  \leanok
  (a) At interior points both notions are the derivative; at the end points both ask for a
  continuous extension. (b) Every point of $U$ is the limit of a sequence in $U^0$, so the
  continuous extensions of the differentials take values in $F_0$ when those on $U^0$ do.
\end{proof}

\section{\texorpdfstring{$\Ck{k,\ell}$}{Ckl} maps and the exponential law}

\begin{definition}[$\Ck{k,\ell}$ maps on sets with dense interior {\cite[Definition~1.7.17]{GN26}}]
  \label{def:lc-ckl}
  \lean{LocallyConvex.Analysis.ContDiffLC₂, LocallyConvex.Analysis.iteratedDLC₂,
    LocallyConvex.Analysis.contDiffLC₂_iff, LocallyConvex.Analysis.contDiffLC₂_iff_contDiffMBOn₂,
    LocallyConvex.Analysis.ContDiffLC₂UpTo, LocallyConvex.Analysis.contDiffLC₂UpTo_coe_iff}
  \leanok
  \uses{def:ckl, def:lc-subset}
  Let $U \subseteq E_1$ and $V \subseteq E_2$ be locally convex subsets with dense interior. A
  map $f \colon U \times V \to F$ is $\Ck{k,\ell}$ if it is continuous, it is $\Ck{k,\ell}$ on
  $U^0 \times V^0$ (Definition~\ref{def:ckl}), and the mixed differentials
  $\dd^{(i,j)}(f|_{U^0 \times V^0})$, $i \le k$, $j \le \ell$, extend to continuous maps on
  $U \times V \times E_1^i \times E_2^j$. For $k, \ell \in \N_0 \cup \{\infty\}$, $f$ is
  $\Ck{k,\ell}$ if it is $\Ck{i,j}$ for all finite $i \le k$, $j \le \ell$. On open sets this is
  Definition~\ref{def:ckl}.
\end{definition}

\begin{theorem}[Properties of $\Ck{k,\ell}$ maps {\cite[Lemma~1.7.18, Remark~1.7.19(b), Propositions~1.7.20, 1.7.26, 1.7.28, 1.7.33(a)]{GN26}}]
  \label{thm:lc-ckl}
  \lean{LocallyConvex.Analysis.ContDiffLC₂.contDiffLC_right,
    LocallyConvex.Analysis.ContDiffLC₂.contDiffLC_left,
    LocallyConvex.Analysis.hasDirDerivAt_of_continuousOn_prod,
    LocallyConvex.Analysis.ContDiffLC₂.contDiffLC₂_partialDLC₁, LocallyConvex.Analysis.ContDiffLC₂.succ_of,
    LocallyConvex.Analysis.ContDiffLC₂.contDiffLC₂_partialDLC₂,
    LocallyConvex.Analysis.ContDiffLC₂.zero_succ_of,
    LocallyConvex.Analysis.contDiffLC₂_pi_iff, LocallyConvex.Analysis.contDiffLC₂_prodMk_iff,
    LocallyConvex.Analysis.ContDiffLC₂.comp_prodMap, LocallyConvex.Analysis.ContDiffLC₂.comp_left,
    LocallyConvex.Analysis.ContDiffLC.contDiffLC₂,
    LocallyConvex.Analysis.ContDiffLC₂UpTo.comp_prodMap,
    LocallyConvex.Analysis.ContDiffLC₂UpTo.comp_left}
  \leanok
  \uses{def:lc-ckl, thm:lc-higher, thm:lc-rules}
  Let $f \colon U \times V \to F$ be $\Ck{k,\ell}$.
  \begin{enumerate}
    \item For every $x \in U$, also on the boundary, $f(x, \cdot)$ is $\Ck{\ell}$ with
      $\dd^{(j)}(f(x, \cdot)) = \dd^{(0,j)} f(x, \cdot)$; for $y \in V$ and $w \in E_2^j$,
      $x \mapsto \dd^{(0,j)} f(x, y, w)$ is $\Ck{k}$.
    \item \textbf{Recursion.} $f$ is $\Ck{k+1,\ell}$ if and only if it is $\Ck{0,\ell}$ and
      $((x, a), y) \mapsto \dd^{(1,0)} f(x, y, a)$ is $\Ck{k,\ell}$ on
      $(U \times E_1) \times V$; and
      similarly in the second variable.
    \item A map into a product is $\Ck{k,\ell}$ if and only if its components are.
    \item \textbf{Chain rules.} $(x, y) \mapsto g(f_1(x), f_2(y))$ is $\Ck{r,s}$ if $g$ is
      $\Ck{r,s}$, $f_1$ is $\Ck{r}$ and $f_2$ is $\Ck{s}$; and $g \circ f$ is $\Ck{r,s}$ if $f$ is
      $\Ck{r,s}$ and $g$ is $\Ck{r+s}$.
    \item $\Ck{k+\ell}$ maps on $U \times V$ are $\Ck{k,\ell}$.
  \end{enumerate}
\end{theorem}

\begin{proof}
  \leanok
  (a) Pass to the limit from $U^0$ with a parametrized form of the mean value inclusion: a family
  of maps with continuous directional derivatives keeps them at accumulation points of the
  parameters (the role of \cite[Lemma~1.4.18]{GN26}); conclude with Theorem~\ref{thm:lc-higher}.
  (b) From explicit formulas for the mixed differentials of the partial differentials, and the
  interchange of a derivative in the first variable with derivatives in the second. (c)--(e)
  follow from (b) by induction, with the chain rule of Theorem~\ref{thm:lc-rules}.
\end{proof}

\begin{remark}
  Glöckner and Neeb prove the chain rules by computing the mixed differentials of the composite
  with Faà di Bruno's formula \cite[(1.91)--(1.97)]{GN26}. The recursive characterization (b)
  avoids those formulas. We do not claim this argument is new.
\end{remark}

\begin{theorem}[Exponential law on sets with dense interior {\cite[Theorem~1.7.30]{GN26}}]
  \label{thm:lc-exponential-law}
  \lean{LocallyConvex.Analysis.ContDiffLCMap, LocallyConvex.Analysis.ContDiffLCMap₂,
    LocallyConvex.Analysis.curryLC, LocallyConvex.Analysis.ContDiffLC₂.contDiffLC_curryLC,
    LocallyConvex.Analysis.iteratedDLC_iteratedDLC_curryLC,
    LocallyConvex.Analysis.ContDiffLCMap₂.isEmbedding_curry,
    LocallyConvex.Analysis.ContDiffLCMap₂.curryEquiv,
    LocallyConvex.Analysis.SmoothLCMap, LocallyConvex.Analysis.SmoothLCMap₂,
    LocallyConvex.Analysis.SmoothLCMap₂.isEmbedding_curry}
  \leanok
  \uses{thm:lc-ckl, thm:exponential-law}
  Give $\Ck{k}(V, F)$ and $\Ck{k,\ell}(U \times V, F)$ the compact-open $\Ck{k}$ and $\Ck{k,\ell}$
  topologies, as in Chapter~\ref{chap:mapping}. For $f \in \Ck{k,\ell}(U \times V, F)$, the map
  $f^\vee \colon U \to \Ck{\ell}(V, F)$, $f^\vee(x) = f(x, \cdot)$, is $\Ck{k}$, with
  \[
    \dd^{(j)}\bigl(\dd^{(i)}(f^\vee)(x, v)\bigr)(y, w) = \dd^{(i,j)} f(x, y, v, w)
    \qquad \text{\cite[(1.98)]{GN26}},
  \]
  and $\Phi \colon f \mapsto f^\vee$ is a linear topological embedding
  $\Ck{k,\ell}(U \times V, F) \to \Ck{k}(U, \Ck{\ell}(V, F))$. If $V$ is locally compact, $\Phi$ is
  an isomorphism of topological vector spaces. The same holds for $k = \ell = \infty$.
\end{theorem}

\begin{proof}
  \leanok
  At interior points the derivative of $f^\vee$ is $(\dd^{(1,0)} f)^\vee$: the difference
  quotients converge in $\Ck{\ell}(V, F)$, uniformly on compact subsets of $V \times E_2^j$, by
  the mean value inclusion with a parameter ranging over all of $V$. The induction then runs
  through the recursion of Theorem~\ref{thm:lc-ckl}(b). For $k = \ell = \infty$, $\Ck{\infty}(V, F)$
  is the closed subspace of compatible families in $\prod_j \Ck{j}(V, F)$, and the law follows from
  the finite orders.
\end{proof}

\begin{remark}
  Glöckner and Neeb differentiate $f^\vee$ through the continuous extension of difference
  quotients of $\Ck{1,0}$ maps \cite[Lemma~1.7.6]{GN26} and an induction on the order of the
  differentials \cite[(1.99)]{GN26}. Their alternative hypothesis (b) on $k_{\R}$-spaces is also
  formalized (\texttt{curryEquivOfKRSpace}); first countable spaces are $k_{\R}$-spaces, so the law
  holds for Fréchet spaces.
\end{remark}
